% ---------------------------------------------------------------
% TCC1 - SENAC SANTO AMARO
% Versão adaptada para o projeto do Matheus
% Tema: classificação de câncer de mama em mamografias
% ---------------------------------------------------------------

\documentclass[
    12pt,
    openright,
    oneside,
    a4paper,
    english,
    brazil
]{abntex2}

% ---------------------------------------------------------------
% PACOTES
% ---------------------------------------------------------------
\usepackage{lmodern}
\usepackage[T1]{fontenc}
\usepackage[utf8]{inputenc}
\usepackage{lastpage}
\usepackage{indentfirst}
\usepackage{color}
\usepackage{graphicx}
\usepackage{microtype}
\usepackage{float}
\usepackage{booktabs}
\usepackage{multirow}
\usepackage{longtable}
\usepackage{amsmath,amssymb}
\usepackage{pgfgantt}
\usepackage[brazilian,hyperpageref]{backref}
\usepackage[alf]{abntex2cite}
\usepackage[margin=3cm]{geometry}

% ---------------------------------------------------------------
% CONFIGURAÇÕES DO PDF
% ---------------------------------------------------------------
\makeatletter
\hypersetup{
    pdftitle={\@title},
    pdfauthor={\@author},
    pdfsubject={TCC1 - SENAC Santo Amaro},
    pdfcreator={LaTeX with abnTeX2},
    pdfkeywords={tcc}{senac}{câncer de mama}{cnn}{computação quântica},
    colorlinks=true,
    linkcolor=blue,
    citecolor=blue,
    filecolor=magenta,
    urlcolor=blue,
    bookmarksdepth=4
}
\makeatother

\definecolor{blue}{RGB}{41,5,195}

% ---------------------------------------------------------------
% INFORMAÇÕES DO DOCUMENTO
% ---------------------------------------------------------------
\titulo{Classificação de câncer de mama em mamografias: um estudo comparativo entre CNNs clássicas e híbridas quântico-clássicas}
\autor{Matheus Sanchez Duda}
\local{São Paulo}
\data{2026}
\orientador{Prof. [Nome do Orientador]}
% \coorientador{Prof. [Nome do Coorientador]}

\instituicao{%
  SENAC -- Serviço Nacional de Aprendizagem Comercial
  \par
  Unidade Santo Amaro
  \par
  Bacharelado em Ciência da Computação
}

\tipotrabalho{Trabalho de Conclusão de Curso I}

\preambulo{Proposta de Trabalho de Conclusão de Curso apresentada ao SENAC Santo Amaro como requisito parcial para aprovação na disciplina de TCC1 do curso de Ciência da Computação.}

\makeindex

\begin{document}
\selectlanguage{brazil}
\frenchspacing

% ---------------------------------------------------------------
% ELEMENTOS PRÉ-TEXTUAIS
% ---------------------------------------------------------------
\pretextual
\imprimircapa
\imprimirfolhaderosto

\begin{resumo}

Este trabalho tem o objetivo de aprimorar a arquitetura de redes neurais convolucionais feita anteriormente, que conseguia classificar imagens de cancer de mama em três classes (normal, benigna, maligna) com uma precisao de cerca de 70 por cento na media para as três classes. Agora deve ser feito a melhoria nos dois tipos de modelos para o classico conseguir clasificar mais exames radiograficos, e do modelo quântico poder ter resultados além do chute aleatório permitindo uma comparação mais robusta entre os dois tipos de modelos. 
\vspace{\onelineskip}
\noindent
\textbf{Palavras-chave}: câncer de mama. mamografia. redes neurais convolucionais. computação quântica. aprendizado de máquina quântico.
\end{resumo}

\pdfbookmark[0]{\contentsname}{toc}
\tableofcontents*
\cleardoublepage

% ---------------------------------------------------------------
% ELEMENTOS TEXTUAIS
% ---------------------------------------------------------------
\textual

% ===============================================================
\chapter{Introdução}
\label{cap:introducao}

Este trabalho propõe a continuidade e o aprofundamento de uma iniciação científica dedicada à classificação de câncer de mama em imagens de mamografia, com comparação entre modelos de redes neurais convolucionais clássicas e modelos híbridos quântico-clássicos. Na implementação atual da pesquisa, a classificação está organizada em três classes principais; no TCC, pretende-se manter esse ponto de partida e investigar, de forma incremental, a ampliação para um número maior de classes. Nesta proposta inicial, são apresentados o contexto do problema, a justificativa do estudo e os objetivos que orientarão o desenvolvimento do TCC.

\section{Contexto}
\label{sec:contexto}

O câncer de mama permanece como um dos principais desafios contemporâneos em saúde pública. A Organização Mundial da Saúde destaca que a doença ocorre em todos os países e que a detecção precoce, associada ao tratamento oportuno, é essencial para reduzir a mortalidade e melhorar os desfechos clínicos \cite{WHO2025BreastCancer}. No contexto brasileiro, o Instituto Nacional de Câncer mantém a mamografia como componente central das estratégias de rastreamento populacional, reforçando a importância do exame para identificação de lesões em estágios iniciais \cite{INCA2025Rastreamento}.

Paralelamente, a área de Inteligência Artificial tem ampliado sua presença na análise de imagens médicas, sobretudo com o uso de redes neurais convolucionais, que se destacam na extração automática de padrões visuais complexos \cite{LeCun2015}. Em mamografias, tais modelos vêm sendo estudados para apoio à classificação de exames, triagem de casos suspeitos e redução de variabilidade entre análises. Entretanto, a literatura recente também aponta limitações importantes relacionadas à reprodutibilidade, à generalização e à interpretação clínica desses sistemas \cite{Bhalla2023MammographyReview}.

Nos últimos anos, a computação quântica passou a ser investigada como alternativa complementar para problemas de aprendizado de máquina, originando modelos híbridos que combinam camadas clássicas e quânticas em uma mesma arquitetura. Apesar do interesse crescente, revisões recentes indicam que ainda existem poucas evidências robustas de superioridade prática dos modelos quânticos em cenários realistas de saúde digital, especialmente quando se considera escalabilidade, ruído e custo computacional \cite{Gupta2025QMLDigitalHealth}. Esse cenário delimita o problema deste TCC: compreender, em uma tarefa específica de classificação de mamografias, se arquiteturas híbridas quântico-clássicas oferecem vantagens reais em comparação com CNNs clássicas, inclusive em um possível cenário de aumento da granularidade das classes.

\section{Justificativa}
\label{sec:justificativa}

A justificativa deste trabalho pode ser compreendida em três dimensões. A primeira é a relevância social do tema. O câncer de mama possui alta incidência e depende fortemente de estratégias de detecção precoce e apoio à decisão clínica, o que torna pertinente investigar métodos computacionais capazes de auxiliar a análise de mamografias \cite{WHO2025BreastCancer}. Ainda que o presente trabalho não tenha objetivo de uso clínico imediato, ele se insere em uma linha de pesquisa com impacto potencial na área de saúde.

A segunda dimensão é a relevância acadêmica. O aprendizado profundo consolidou-se como uma abordagem central em visão computacional, mas a adoção de modelos híbridos quântico-clássicos ainda carece de validação comparativa rigorosa, principalmente em aplicações médicas \cite{LeCun2015,Gupta2025QMLDigitalHealth}. Assim, o TCC contribui ao investigar se a integração de componentes quânticos traz benefícios mensuráveis em um problema concreto e de alta relevância.

A terceira dimensão é a lacuna científica identificada. Estudos sobre deep learning em mamografia apontam desempenho promissor, porém também evidenciam limitações de reprodutibilidade e de avaliação em condições realistas \cite{Bhalla2023MammographyReview}. Na direção quântica, trabalhos recentes mostram resultados competitivos em tarefas de classificação de imagens, mas ainda em cenários restritos, com bases menores, forte dependência de simulação e poucas evidências em contextos próximos do uso aplicado \cite{Long2025QCQCNN}. Dessa forma, este TCC busca preencher uma lacuna ao propor uma comparação experimental controlada, reprodutível e metodologicamente explícita entre abordagens clássicas e híbridas no domínio da mamografia.

\section{Objetivo}
\label{sec:objetivo}


\subsection{Objetivo principal}
\label{subsec:objetivo_principal}

Investigar comparativamente o desempenho de CNNs clássicas e modelos híbridos quântico-clássicos na classificação de câncer de mama em imagens de mamografia, considerando métricas de desempenho preditivo, custo computacional, estabilidade experimental e viabilidade prática, com base inicial no cenário atual de três classes e possibilidade de ampliação para uma classificação mais detalhada.

\subsection{Objetivos secundários}
\label{subsec:objetivos_secundarios}

\begin{enumerate}
    \item Aprimorar a arquitetura clássica de CNN para servir como;
    \item Aprimorar a arquitetura híbrida quântico-clássica compatível com o problema estudado;
    \item Estabelecer um protocolo experimental padronizado para treinamento, validação e teste dos modelos;
    \item Investigar a viabilidade de ampliar o problema atual de três classes para uma classificação com maior granularidade, conforme a disponibilidade e a confiabilidade dos rótulos da base;
    \item Comparar os modelos por meio de métricas como acurácia, precisão, recall, F1-score, AUC e matriz de confusão;
    \item Avaliar também o custo computacional, o tempo de treinamento, a estabilidade entre execuções e as limitações de escalabilidade;
    \item Discutir criticamente os resultados obtidos, identificando vantagens, desvantagens e limitações práticas da abordagem híbrida.
\end{enumerate}

% ===============================================================
\chapter{Revisão Bibliográfica}
\label{cap:revisao}


\section{Referencial teórico}
\label{sec:referencial}

\subsection{Câncer de mama e mamografia}

O câncer de mama é uma doença caracterizada pelo crescimento descontrolado de células mamárias, podendo evoluir para formas invasivas e metastáticas. A Organização Mundial da Saúde ressalta que a mortalidade pode ser reduzida com diagnóstico e tratamento precoces, sendo a mamografia um dos principais instrumentos de rastreamento em populações assintomáticas \cite{WHO2025BreastCancer}. No Brasil, o INCA mantém orientações específicas para o rastreamento mamográfico, o que reforça a relevância da mamografia no contexto nacional \cite{INCA2025Rastreamento}.

No campo computacional, imagens de mamografia são particularmente desafiadoras por apresentarem ruído, variações de contraste, estruturas anatômicas complexas e, em muitos casos, desbalanceamento entre classes. Tais características tornam necessário o uso de métodos robustos de extração de atributos e classificação. Além disso, a definição de mais classes diagnósticas aumenta a dificuldade do problema, exigindo maior cuidado na curadoria dos dados e na avaliação do desempenho.

\subsection{Aprendizado profundo e redes neurais convolucionais}

O avanço recente da visão computacional está fortemente associado ao desenvolvimento do aprendizado profundo, especialmente com o uso de redes neurais convolucionais. Conforme discutido por \citeonline{LeCun2015}, redes profundas aprendem representações hierárquicas diretamente a partir dos dados, dispensando parte da engenharia manual de atributos tradicional.

Em aplicações médicas, CNNs têm sido empregadas para classificação, detecção e segmentação de padrões em exames de imagem. Na área de mamografia, revisões sistemáticas indicam amplo uso dessas arquiteturas, porém alertam para dificuldades de reprodutibilidade, viés de seleção e falta de avaliação em ambientes mais próximos da prática clínica \cite{Bhalla2023MammographyReview}. Esses pontos justificam a adoção de um protocolo experimental rigoroso neste TCC.

\subsection{Computação quântica e modelos híbridos}

A computação quântica explora propriedades como superposição, emaranhamento e interferência para processamento de informação. Em aprendizado de máquina, uma das linhas mais promissoras tem sido a construção de modelos híbridos, nos quais etapas clássicas e quânticas são combinadas para aproveitar, ao menos em tese, diferentes capacidades de representação.

Entretanto, a literatura recente sugere cautela. A revisão sistemática de \citeonline{Gupta2025QMLDigitalHealth} mostra que, em saúde digital, a maioria dos estudos ainda opera com fortes simplificações, pequeno porte de dados e escopo experimental reduzido. Já \citeonline{Long2025QCQCNN} apresenta uma arquitetura híbrida quântico-clássico-quântica com resultados competitivos em bases de imagens e tumores por ressonância magnética, evidenciando potencial, mas também dependência de cenários controlados.

Assim, a comparação entre uma CNN clássica e uma arquitetura híbrida em mamografias torna-se um recorte adequado para investigar, de forma mais controlada, o real benefício dessas abordagens.

\section{Metodologia da pesquisa bibliográfica}
\label{sec:metodologia_biblio}

A pesquisa bibliográfica será conduzida em bases como Google Scholar, PubMed, IEEE Xplore, ACM Digital Library, Scopus, SciELO e Periódicos CAPES. Serão utilizadas palavras-chave em português e inglês, incluindo expressões como \textit{``câncer de mama''}, \textit{``mamografia''}, \textit{``breast cancer''}, \textit{``mammography''}, \textit{``convolutional neural network''}, \textit{``deep learning''}, \textit{``quantum machine learning''}, \textit{``hybrid quantum-classical neural network''} e combinações derivadas.

Como critérios de inclusão, serão priorizados trabalhos publicados entre 2018 e 2026, em português ou inglês, que abordem diretamente classificação de câncer de mama, mamografia, visão computacional médica ou modelos híbridos quântico-clássicos aplicados à classificação de imagens. Também serão considerados estudos sobre bases de dados amplas e cenários multiclasse quando contribuírem para a futura expansão do problema atualmente tratado com três classes. Serão excluídos trabalhos sem relação direta com o problema de pesquisa, textos opinativos sem descrição metodológica mínima e estudos sem resultados experimentais suficientemente claros.

A estratégia de revisão buscará não apenas reunir referências, mas identificar convergências, divergências, lacunas e limitações metodológicas relevantes para posicionar adequadamente a proposta deste TCC.

\section{Fundamentos}
\label{sec:fundamentos}

Do ponto de vista técnico, a solução proposta envolve quatro blocos principais. O primeiro é a preparação dos dados, incluindo padronização das imagens, redimensionamento, normalização e definição dos conjuntos de treino, validação e teste. Nessa etapa, será preservado inicialmente o cenário atual de três classes, com estudo posterior de expansão para mais classes caso a base escolhida ofereça rótulos consistentes para isso. O segundo é a construção do modelo clássico de referência, provavelmente baseado em uma CNN com número controlado de parâmetros e treinamento supervisionado.

O terceiro bloco consiste na definição do modelo híbrido quântico-clássico. Nessa etapa, serão estudadas alternativas como camadas variacionais quânticas, estratégias de codificação de dados em poucos qubits e integração com camadas convolucionais ou densas clássicas. O objetivo não é propor uma arquitetura clínica final, mas construir um modelo comparável ao \textit{baseline} clássico dentro das restrições de hardware e tempo disponíveis.

O quarto bloco compreende o protocolo de avaliação. Serão empregadas métricas de classificação como acurácia, precisão, recall, F1-score, AUC e matriz de confusão, além de medidas relacionadas a tempo de treinamento, tempo de inferência, estabilidade entre execuções e custo computacional. Também se pretende registrar o ambiente experimental, versões de bibliotecas e sementes aleatórias para favorecer a reprodutibilidade.



% ===============================================================
\chapter{Desenvolvimento}
\label{cap:desenvolvimento}


\section{Solução proposta}
\label{sec:solucao}


\section{Cronograma de trabalho}
\label{sec:cronograma}

% ---------------------------------------------------------------
% ELEMENTOS PÓS-TEXTUAIS
% ---------------------------------------------------------------
\postextual
\bibliography{bibliografia_matheus_v2}

\end{document}
